\begin{table*}[t]\centering\caption{Primary regime, 1:20 (12 training positives), 10 seeds. AUROC, AUPRC, Brier score and $|a|$ (absolute calibration-in-the-large): mean $\pm$ sd; calibration slope $b$ and $|b-1|$: median [IQR]. Threshold moving has the probabilities of no correction and is omitted.}\label{tab:main}
\small\begin{tabular}{lcccccc}
\toprule
System & AUROC & AUPRC & Brier & slope $b$ & $|b-1|$ & $|a|$ \\
\midrule
LR + no correction & 0.991 $\pm$ 0.005 & 0.941 $\pm$ 0.030 & 0.008 $\pm$ 0.003 & 1.10 [0.99, 1.23] & 0.10 [0.04, 0.23] & 0.748 $\pm$ 0.349 \\
LR + reweight & 0.989 $\pm$ 0.007 & 0.934 $\pm$ 0.034 & 0.010 $\pm$ 0.003 & 0.90 [0.77, 1.02] & 0.22 [0.07, 0.27] & 0.805 $\pm$ 0.520 \\
LR + ROS & 0.989 $\pm$ 0.009 & 0.933 $\pm$ 0.033 & 0.009 $\pm$ 0.004 & 0.81 [0.67, 0.90] & 0.27 [0.10, 0.35] & 0.663 $\pm$ 0.506 \\
LR + SMOTE & 0.990 $\pm$ 0.007 & 0.939 $\pm$ 0.032 & 0.008 $\pm$ 0.003 & 0.82 [0.70, 0.91] & 0.28 [0.16, 0.34] & 0.565 $\pm$ 0.388 \\
\midrule
GB + no correction & 0.909 $\pm$ 0.057 & 0.732 $\pm$ 0.170 & 0.022 $\pm$ 0.009 & 0.32 [0.25, 0.43] & 0.68 [0.57, 0.75] & 5.114 $\pm$ 3.414 \\
GB + reweight & 0.890 $\pm$ 0.087 & 0.757 $\pm$ 0.156 & 0.023 $\pm$ 0.013 & 0.39 [0.31, 0.51] & 0.61 [0.49, 0.69] & 3.877 $\pm$ 2.772 \\
GB + ROS & 0.914 $\pm$ 0.050 & 0.764 $\pm$ 0.107 & 0.025 $\pm$ 0.013 & 0.36 [0.29, 0.47] & 0.64 [0.53, 0.71] & 4.706 $\pm$ 2.827 \\
GB + SMOTE & 0.911 $\pm$ 0.048 & 0.734 $\pm$ 0.154 & 0.029 $\pm$ 0.014 & 0.31 [0.25, 0.43] & 0.69 [0.57, 0.75] & 5.238 $\pm$ 1.467 \\
\bottomrule
\end{tabular}
\end{table*}

\begin{table*}[t]\centering\caption{AUROC, AUPRC and Brier score across prevalences (nat.\ = natural), mean of 10 seeds; threshold moving equals no correction here.}\label{tab:regimes}
\small\begin{tabular}{lcccccccccccc}
\toprule
 & \multicolumn{4}{c}{AUROC} & \multicolumn{4}{c}{AUPRC} & \multicolumn{4}{c}{Brier} \\
System & nat. & 1:5 & 1:20 & 1:50 & nat. & 1:5 & 1:20 & 1:50 & nat. & 1:5 & 1:20 & 1:50 \\
\midrule
LR + no correction & 0.993 & 0.993 & 0.991 & 0.989 & 0.992 & 0.980 & 0.941 & 0.910 & 0.024 & 0.016 & 0.008 & 0.005 \\
LR + reweight & 0.993 & 0.991 & 0.989 & 0.988 & 0.992 & 0.976 & 0.934 & 0.899 & 0.025 & 0.019 & 0.010 & 0.006 \\
LR + ROS & 0.993 & 0.990 & 0.989 & 0.988 & 0.992 & 0.974 & 0.933 & 0.895 & 0.025 & 0.020 & 0.009 & 0.006 \\
LR + SMOTE & 0.993 & 0.991 & 0.990 & 0.990 & 0.992 & 0.976 & 0.939 & 0.910 & 0.025 & 0.018 & 0.008 & 0.005 \\
\midrule
GB + no correction & 0.986 & 0.977 & 0.909 & 0.878 & 0.984 & 0.941 & 0.732 & 0.638 & 0.042 & 0.040 & 0.022 & 0.013 \\
GB + reweight & 0.986 & 0.977 & 0.890 & 0.849 & 0.983 & 0.943 & 0.757 & 0.570 & 0.041 & 0.035 & 0.023 & 0.015 \\
GB + ROS & 0.986 & 0.981 & 0.914 & 0.834 & 0.984 & 0.951 & 0.764 & 0.562 & 0.041 & 0.036 & 0.025 & 0.015 \\
GB + SMOTE & 0.987 & 0.981 & 0.911 & 0.850 & 0.985 & 0.950 & 0.734 & 0.553 & 0.039 & 0.032 & 0.029 & 0.016 \\
\bottomrule
\end{tabular}
\end{table*}

\begin{table*}[t]\centering\caption{Paired difference to no correction, same learner (positive $\Delta$Brier = worse). $^{*}$: unadjusted paired $p<0.05$, 10 seeds.}\label{tab:deltas}
\small\begin{tabular}{lccccccccc}
\toprule
 & \multicolumn{3}{c}{1:5} & \multicolumn{3}{c}{1:20} & \multicolumn{3}{c}{1:50} \\
Correction & $\Delta$AUROC & $\Delta$AUPRC & $\Delta$Brier & $\Delta$AUROC & $\Delta$AUPRC & $\Delta$Brier & $\Delta$AUROC & $\Delta$AUPRC & $\Delta$Brier \\
\midrule
LR + reweight & -0.002\rlap{$^{*}$} & -0.004\rlap{$^{*}$} & +0.003\rlap{$^{*}$} & -0.002 & -0.007 & +0.001 & -0.001 & -0.011 & +0.001 \\
LR + ROS & -0.003\rlap{$^{*}$} & -0.006\rlap{$^{*}$} & +0.004\rlap{$^{*}$} & -0.002 & -0.008 & +0.001 & -0.001 & -0.015 & +0.001 \\
LR + SMOTE & -0.001 & -0.004 & +0.002 & -0.001 & -0.002 & -0.000 & +0.001 & -0.000 & +0.000 \\
GB + reweight & +0.000 & +0.002 & -0.004 & -0.019 & +0.025 & +0.002 & -0.029 & -0.068 & +0.002 \\
GB + ROS & +0.004\rlap{$^{*}$} & +0.011\rlap{$^{*}$} & -0.003 & +0.005 & +0.031 & +0.003 & -0.044 & -0.076 & +0.002 \\
GB + SMOTE & +0.004 & +0.010 & -0.007 & +0.002 & +0.002 & +0.007 & -0.028 & -0.085 & +0.003 \\
\bottomrule
\end{tabular}
\end{table*}

\textbf{H1 (corrections do not improve discrimination): mostly supported for LR, mixed for GB.} For LR the AUROC differences to no correction are at most 0.003 in absolute value (Table~\ref{tab:deltas}; at natural prevalence all LR systems have AUROC 0.993, Table~\ref{tab:regimes}), and at 1:5 reweighting and ROS are slightly worse in AUROC, AUPRC and Brier score ($p\approx0.03$--$0.05$ unadjusted). For GB, 11 of the 18 AUROC/AUPRC differences at 1:5--1:50 are positive. The registered refutation criterion (mean paired gain $\ge 0.01$ with $p<0.05$) is met by a single cell, GB with ROS at 1:5 (AUPRC $+0.011$), one of 18 unadjusted tests on GB, so we cannot separate it from chance. The gains at 1:20 (AUPRC $+0.025$ reweight, $+0.031$ ROS) are not significant (seed sd of AUPRC 0.11--0.17, Table~\ref{tab:main}). At 1:50 all three corrections are \emph{worse} than no correction for GB (AUROC $-0.028$ to $-0.044$, AUPRC $-0.068$ to $-0.085$), again not significant ($p>0.2$). No correction reliably improves discrimination here.

\textbf{H2 (corrections worsen Brier score and $|a|$ at 1:20 and 1:50): direction as predicted for Brier but not significant; not supported for $|a|$.} Brier score is higher with a correction in 11 of the 12 cells at 1:20 and 1:50; the exception, which refutes H2 for that cell under the registered rule, is LR with SMOTE at 1:20 ($-0.000$ in Table~\ref{tab:deltas}). Effects are small (LR $+0.001$, GB $+0.002$ to $+0.007$ at 1:20) and no cell has $p<0.05$ (smallest $p=0.085$, LR reweighting at 1:20). Mean $|a|$ is higher with a correction in only 5 of these 12 cells and no difference is significant (Table~\ref{tab:citl}, smallest $p=0.11$).

\textbf{Signed calibration-in-the-large (not registered).} $|a|$ hides a clear effect. LR without correction underestimates risk at every prevalence ($a=+0.33$ to $+0.76$, Table~\ref{tab:citl}). Every correction lowers $a$ in all 12 LR cells ($p\le0.0002$): at natural prevalence $a$ moves toward zero ($-0.06$ to $+0.06$), and at 1:5, 1:20 and 1:50 it changes sign ($-0.35$ to $-0.89$), i.e.\ the corrected LR models overestimate risk, the direction reported in~\cite{carriero2024harms,sirikul2026class}. Here an underestimate is replaced by an overestimate of similar size, so $|a|$ and the Brier score barely move. For GB, $a$ is positive in every cell and lower with a correction at natural prevalence ($p\le0.0002$); at 1:5--1:50 only SMOTE at 1:5 differs with $p<0.05$, and $|a|$ is 1.3 to 5.7 with or without correction.

\textbf{Calibration slope.} For LR without correction the median slope is above 1 at every prevalence (Fig.~\ref{fig:slope}) (predictions too moderate); the corrections lower the median below 1 at 1:5--1:50 (at 1:20 from 1.10 to 0.81--0.90, Table~\ref{tab:main}). The signed shift is significant in all nine LR cells at 1:5--1:50 ($p=0.0002$--$0.027$) and absent at natural prevalence ($p>0.5$). The registered distance $|b-1|$ does not change significantly in any LR cell ($p\ge0.16$): at 1:20 its median is larger with a correction (0.10 versus 0.22--0.28) while its mean is smaller (\texttt{deltas\_all.csv}), so we draw no conclusion about which is closer to 1. GB is over-confident with or without correction (median slope 0.31--0.39 at 1:20).

\begin{figure}[t]\centering\includegraphics[width=\columnwidth]{fig_slope.pdf}
\caption{Calibration slope (log axis), median and IQR over 10 seeds by prevalence; dashed line = 1.}\label{fig:slope}\end{figure}

\textbf{H3 (threshold moving): supported, with a trivial first part.} ``Threshold'' and ``no correction'' have identical AUROC, AUPRC, Brier score and slope in every run (\texttt{h3\_maxabsdiff.csv}); both arms use the same fitted model, so this verifies the implementation rather than testing a hypothesis. Moving the cut-off to the training prevalence raises sensitivity and lowers specificity (Table~\ref{tab:oper}). For LR sensitivity rises from 0.727 to 0.909 at 1:50 (specificity 1.000 to 0.967) and balanced accuracy from 0.863 to 0.938; at natural prevalence the balanced-accuracy change is within noise (0.963 to 0.965, $p=0.46$). For GB the gains are smaller and at 1:20 and 1:50 not significant ($p\approx0.07$--$0.10$).

\textbf{Operating point of the corrections (not registered).} Threshold moving is not the only strategy that moves the LR operating point. At the 0.5 cut-off reweighting, ROS and SMOTE raise LR sensitivity in all 12 cells ($p\le0.012$; 0.727 to 0.812--0.816 at 1:50) and LR balanced accuracy at 1:5, 1:20 and 1:50 ($p=0.0003$--$0.032$; 0.863 to 0.905--0.906 at 1:50), with a specificity loss that is significant in 7 of the 12 cells. Threshold moving raises both further: balanced accuracy is 0.947 against 0.927--0.928 for the corrections at 1:20 and 0.938 against 0.905--0.906 at 1:50 (threshold versus each correction $p<0.01$ at 1:20 and 1:50, $p=0.014$--$0.077$ at 1:5, no difference at natural prevalence), with lower specificity than the corrections ($p<0.05$ in all 12 cells) and without changing a probability. This agrees with the signed $a$: the corrections act on LR mainly as an upward shift of risk. For GB the corrections do not help: sensitivity and balanced accuracy improve with $p<0.05$ only for SMOTE at 1:5, and at 1:20 and 1:50 all three corrections have lower mean balanced accuracy than no correction ($p>0.2$).

\begin{table*}[t]\centering\caption{Operating point at each system's cut-off (0.5; training prevalence for threshold moving), mean of 10 seeds. $^{*}$: unadjusted paired $p<0.05$ against no correction, same learner.}\label{tab:oper}
\small\begin{tabular}{lcccccccccccc}
\toprule
 & \multicolumn{4}{c}{Sensitivity} & \multicolumn{4}{c}{Specificity} & \multicolumn{4}{c}{Balanced accuracy} \\
System & nat. & 1:5 & 1:20 & 1:50 & nat. & 1:5 & 1:20 & 1:50 & nat. & 1:5 & 1:20 & 1:50 \\
\midrule
LR + no correction & 0.939 & 0.887 & 0.805 & 0.727 & 0.988 & 0.997 & 1.000 & 1.000 & 0.963 & 0.942 & 0.902 & 0.863 \\
LR + reweight & 0.955\rlap{$^{*}$} & 0.928\rlap{$^{*}$} & 0.861\rlap{$^{*}$} & 0.816\rlap{$^{*}$} & 0.980\rlap{$^{*}$} & 0.985\rlap{$^{*}$} & 0.993\rlap{$^{*}$} & 0.997 & 0.968\rlap{$^{*}$} & 0.957\rlap{$^{*}$} & 0.927\rlap{$^{*}$} & 0.906\rlap{$^{*}$} \\
LR + ROS & 0.950\rlap{$^{*}$} & 0.925\rlap{$^{*}$} & 0.859\rlap{$^{*}$} & 0.812\rlap{$^{*}$} & 0.978\rlap{$^{*}$} & 0.984\rlap{$^{*}$} & 0.995 & 0.997 & 0.964 & 0.955\rlap{$^{*}$} & 0.927\rlap{$^{*}$} & 0.905\rlap{$^{*}$} \\
LR + SMOTE & 0.950\rlap{$^{*}$} & 0.920\rlap{$^{*}$} & 0.858\rlap{$^{*}$} & 0.814\rlap{$^{*}$} & 0.980\rlap{$^{*}$} & 0.987\rlap{$^{*}$} & 0.998 & 0.997 & 0.965 & 0.954\rlap{$^{*}$} & 0.928\rlap{$^{*}$} & 0.906\rlap{$^{*}$} \\
LR + threshold & 0.956\rlap{$^{*}$} & 0.952\rlap{$^{*}$} & 0.922\rlap{$^{*}$} & 0.909\rlap{$^{*}$} & 0.974\rlap{$^{*}$} & 0.975\rlap{$^{*}$} & 0.972\rlap{$^{*}$} & 0.967\rlap{$^{*}$} & 0.965 & 0.963\rlap{$^{*}$} & 0.947\rlap{$^{*}$} & 0.938\rlap{$^{*}$} \\
\midrule
GB + no correction & 0.912 & 0.839 & 0.727 & 0.616 & 0.969 & 0.980 & 0.990 & 0.993 & 0.941 & 0.910 & 0.858 & 0.805 \\
GB + reweight & 0.922 & 0.858 & 0.688 & 0.614 & 0.963 & 0.979 & 0.990 & 0.992 & 0.942 & 0.918 & 0.839 & 0.803 \\
GB + ROS & 0.920 & 0.852 & 0.672 & 0.573 & 0.967 & 0.980 & 0.989 & 0.992 & 0.944 & 0.916 & 0.830 & 0.783 \\
GB + SMOTE & 0.919 & 0.875\rlap{$^{*}$} & 0.716 & 0.572 & 0.964 & 0.980 & 0.981 & 0.991 & 0.942 & 0.928\rlap{$^{*}$} & 0.848 & 0.781 \\
GB + threshold & 0.923\rlap{$^{*}$} & 0.856\rlap{$^{*}$} & 0.762 & 0.652 & 0.964 & 0.974\rlap{$^{*}$} & 0.986 & 0.991 & 0.944\rlap{$^{*}$} & 0.915\rlap{$^{*}$} & 0.874 & 0.821 \\
\bottomrule
\end{tabular}
\end{table*}

\textbf{Where the methods fail.} GB at 1:50 is the clearest failure: with 5 training positives, GB without correction reaches AUROC 0.878 and with any correction 0.834--0.850 (Table~\ref{tab:regimes}), its AUPRC is 0.638 without correction against 0.910 for LR, and its sensitivity at 0.5 is 0.572--0.616 (Table~\ref{tab:oper}). Resampling five positives may make GB overfit them; we did not isolate this. On some seeds the three corrections give identical GB test metrics at 1:50, and the largest slopes in the registry come from this regime.
